\chapter{Topological Structures and Their Logical Languages}
\label{sec:p03-topological-model-theory}

A topology on the carrier of a structure is different from a topology on
the space of codes of ordinary structures. Here the points of the structure
itself have neighborhoods, and formulas can quantify over those
neighborhoods. Relations and operations are not automatically continuous
or closed merely because the carrier has a topology.
The correspondence-language presentation below uses
\cite[\S2.3, pp.~7--12]{TenCateGabelaiaSustretov2006TopologyDraft};
the arbitrary-signature reduced-ideal-product presentation uses
\cite[\S\S1--2]{GrulovicKurilic2001TopologicalFV}.
The first is the separately identified October~2006 author draft, not its
2009 journal version. Neither is identified with Ziegler's original chapter
or Sgro's interior-operator logic.

\section{Point sorts, open sorts, and restricted quantification}

\begin{hbdefinition}{def:p03-topological-carrier}{Topological and basoid models}
In the correspondence-language signature, let \(X\ne\varnothing\),
let \(\tau\subseteq\mathcal P(X)\) be a topology, and let
\(\mathrm{Prop}=\{p_1,p_2,\ldots\}\) be the source's countably
infinite set of propositional letters. Let
\(\nu:\mathrm{Prop}\to\mathcal P(X)\) assign subsets to the
propositional letters. A topological model is \(M=(X,\tau,\nu)\).
No openness condition is imposed on the sets \(\nu(p)\).
A basis \(\sigma\subseteq\mathcal P(X)\) in this presentation contains
\(\varnothing\), covers \(X\), and has the refinement property:
if \(x\in U\cap V\) for \(U,V\in\sigma\), some
\(W\in\sigma\) contains \(x\) and is contained in \(U\cap V\).
Its arbitrary unions form the generated topology \(\widehat\sigma\).
A basoid model is \((X,\sigma,\nu)\); write
\(\widehat M=(X,\widehat\sigma,\nu)\).
Requiring the empty basis element is a convention of this presentation;
adding it to an ordinary basis does not change its topology.
\end{hbdefinition}

\begin{hbdefinition}{def:p03-topological-languages}{The two-sorted language and membership polarity}
The language \(L_2\) has point variables \(x,y,z\), open variables
\(U,V,W\), unary predicates \(P_p\), point equality, open equality,
and membership \(x\varepsilon U\). It has ordinary finitary classical
connectives and quantifiers in each sort. In a topological model, point
quantifiers range over \(X\) and open quantifiers over \(\tau\);
in a basoid model the latter range over \(\sigma\).
Membership is actual membership and each equality is actual identity in
its own sort. Interpret \(P_p(x)\) by \(x\in\nu(p)\).

For the restricted language \(L_t\), test polarity on membership atoms:
first expand open equality \(U=V\) as
\(\forall z(z\varepsilon U\leftrightarrow z\varepsilon V)\),
with a fresh point variable, and expand abbreviations for connectives.
Then \(U\) is positive or negative in a formula when all its free
membership occurrences lie under respectively even or odd numbers of
negations. An absent variable has both polarities. Open equality must
not be counted as a monotone membership atom.
The membership presentation of \(L_t\) starts from point equality,
\(P_p(x)\), membership, and truth, and is closed under classical
connectives, point quantification, and these guarded open quantifiers:
\[
 \exists U\,(x\varepsilon U\land\alpha)
       \quad\text{if \(U\) is negative in \(\alpha\)},
 \qquad
 \forall U\,(x\varepsilon U\to\alpha)
       \quad\text{if \(U\) is positive in \(\alpha\)}.
\]
Free open equality is available through its expanded formula; a guarded
formula is admitted only if the expanded polarity test succeeds.
The universal clause is the classical dual of the existential one.
Bound variables are renamed before substitutions.
\end{hbdefinition}

\begin{hblemma}{thm:p03-topological-polarity}{proved=annex:p03-topological-polarity}{Monotonicity in a free open parameter}
In either a topology or a basis, with all other parameters fixed, truth of
a formula positive in \(U\) is preserved when its assigned subset is
enlarged. Truth of a formula negative in \(U\) is preserved when it is
shrunk. In this statement a free open parameter may be any fixed subset of
the carrier; only bound open quantifiers use the stated topology or basis.
Thus the same induction also compares a basis element with a larger open
set that is not itself a basis element.
\end{hblemma}

\begin{hbtheorem}{thm:p03-topological-basis}{proved=annex:p03-topological-basis}{Basis reconstruction and basis invariance}
The unions of a basis as above form a topology. For every \(L_t\)-formula,
the basoid model and its generated topological model agree at identical
point parameters and identical free open parameters. In particular they
agree on all \(L_t\)-sentences. The statement permits arbitrary subset
values for the free parameters during its inductive proof; it therefore
implies the source's assertion for free basis elements.
The whole \(L_2\) is not basis invariant.
\end{hbtheorem}

\begin{hbdefinition}{def:p03-topological-basis-axioms}{First-order axioms for a basis presentation}
The following are two-sorted \(L_2\)-sentences, not assertions that the
open sort already contains every union:
\[
\begin{aligned}
 \mathrm{Ext}&:\quad
 \forall U\forall V\bigl(U=V\leftrightarrow
                  \forall x(x\varepsilon U\leftrightarrow x\varepsilon V)\bigr),\\
 \mathrm{Union}&:\quad\forall x\exists U\,(x\varepsilon U),\\
 \mathrm{Empty}&:\quad\exists U\forall x\,\neg(x\varepsilon U),\\
 \mathrm{Bas}&:\quad
 \forall x\forall U\forall V\bigl((x\varepsilon U\land x\varepsilon V)
       \to\exists W(x\varepsilon W\land
         \forall z(z\varepsilon W\to(z\varepsilon U\land z\varepsilon V)))\bigr).
\end{aligned}
\]
These retain the source's names for its four basis axioms.
\end{hbdefinition}

\begin{hbtheorem}{thm:p03-topological-compactness-ls}{proved=annex:p03-topological-compactness-ls}{Compactness and countable-basis models}
Work in a metatheory with Choice and ordinary first-order compactness and
downward L\"owenheim--Skolem. Every finitely satisfiable set of
\(L_t\)-sentences has a topological model. For an at most countable
signature, a satisfiable set of such sentences has a model with countably
many points and a countable basis; if it has an infinite-point model it
has one with countably infinitely many points.
The topology itself need not be countable: it contains all unions of that
countable basis. The theorem is not compactness or a downward cardinal
bound for unrestricted \(L_2\) on full topologies.
\end{hbtheorem}

\begin{hbtheorem}{thm:p03-topological-property-formulas}{proved=annex:p03-topological-property-formulas}{Separation, discreteness, and least neighborhoods}
Write \(\exists U_x\alpha\) for
\(\exists U(x\varepsilon U\land\alpha)\), and
\(\forall U_x\alpha\) for its guarded universal counterpart.
The following \(L_t\)-sentences express the stated properties:
\[
\begin{aligned}
 T_0 &: \forall x\forall y\bigl(x\ne y\to
           (\exists U_x\neg(y\varepsilon U)\lor
            \exists V_y\neg(x\varepsilon V))\bigr),\\
 T_1 &: \forall x\forall y\bigl(x\ne y\to
                                  \exists U_x\neg(y\varepsilon U)\bigr),\\
 T_2 &: \forall x\forall y\bigl(x\ne y\to
        \exists U_x\exists V_y\forall z
                           (\neg(z\varepsilon U)\lor\neg(z\varepsilon V))\bigr),\\
 \mathrm{Discrete} &: \forall x\exists U_x\forall y
                                     (y\varepsilon U\to y=x),\\
 \mathrm{Alexandroff} &: \forall x\exists U_x\forall V_x\forall y
                                     (y\varepsilon U\to y\varepsilon V).
\end{aligned}
\]
The last property says that arbitrary intersections of opens are open,
equivalently that every point has a least open neighborhood. Its inclusion
direction is \(U\subseteq V\). Reversing that direction would allow the
whole space as \(U\) in every topology and would not define this property.
Unrestricted \(L_2\) can also express connectedness by
\[
 \forall U\forall V\left[
   \forall x(x\varepsilon U\leftrightarrow\neg(x\varepsilon V))
    \to\bigl(\forall x\,x\varepsilon U\lor
              \forall x\,\neg(x\varepsilon U)\bigr)\right].
\]
This latter formula uses unguarded open quantifiers; its displayed form
is not an \(L_t\)-formula.
\end{hbtheorem}

\begin{hbdefinition}{def:p03-topological-modal-translation}{Interior semantics and the source's standard translation}
In the basic modal language \(\mathrm{ML}\), interpret truth by all
points, \(p\) by \(\nu(p)\), conjunction by intersection, negation
by complement, and \(\Box\varphi\) by the interior of the truth set of
\(\varphi\). The translation \(\mathrm{ST}_x\) into \(L_t\) is
\[
\begin{aligned}
 \mathrm{ST}_x(\top)&=\top,&
 \mathrm{ST}_x(p)&=P_p(x),\\
 \mathrm{ST}_x(\neg\varphi)&=\neg\mathrm{ST}_x(\varphi),&
 \mathrm{ST}_x(\varphi\land\psi)&=
                \mathrm{ST}_x(\varphi)\land\mathrm{ST}_x(\psi),\\
 \mathrm{ST}_x(\Box\varphi)&=
       \exists U\bigl(x\varepsilon U\land
              \forall y(y\varepsilon U\to\mathrm{ST}_y(\varphi))\bigr).
\end{aligned}
\]
All newly used variables are fresh. The matrix is negative in \(U\),
so this is an admitted guarded quantifier, not unrestricted second-order
quantification. This is the correspondence translation of
\cite[Definition~14 and Theorem~2.10, p.~12]{TenCateGabelaiaSustretov2006TopologyDraft}.
\end{hbdefinition}

\begin{hbtheorem}{thm:p03-topological-translation-homeomorphism}{proved=annex:p03-topological-translation-homeomorphism}{Translation and homeomorphism invariance}
For every modal formula, model and point,
\[
                   M,a\models\varphi
                 \quad\Longleftrightarrow\quad
                   M\models\mathrm{ST}_x(\varphi)[a].
\]
A bijection of carriers that is a homeomorphism and carries every predicate
interpretation to its matching interpretation preserves and reflects all
\(L_2\)-formulas, hence all \(L_t\)-formulas, with point and open
parameters transported along it. No such full preservation is asserted for
arbitrary continuous maps or merely ordinary first-order isomorphisms.
\end{hbtheorem}

\section{Reduced ideal-products of topological structures}
\label{sec:p03-topological-reduced-products}

This section retains the separate presentation of
\cite[pp.~1--7]{GrulovicKurilic2001TopologicalFV}. Its point signature
\(L\) may have arbitrary finite-arity functions and relations. The point
atoms in \(L_t\) now include all corresponding ordinary first-order
atoms, and guards may use point terms \(t\), not just variables.
Their interpretations do not involve the open parameter, so the polarity,
basis-invariance and whole-carrier substitution proofs apply unchanged.
Continuity of the point operations is not an implicit hypothesis.

\begin{hbdefinition}{def:p03-topological-rip}{Ideal-supported boxes and their filter quotient}
Let \(I\ne\varnothing\) be a set of indices, \(\Lambda\) an ideal
of \(\mathcal P(I)\), allowed to be the whole algebra, and
\(\Psi\) a proper filter of \(\mathcal P(I)\). Every
\((S_i,O_i)\) is a nonempty topological \(L\)-structure. Work with
Choice. A basic \(\Lambda\)-box is
\[
       B=\prod_i U_i,\qquad U_i\in O_i,\qquad
                \{i:U_i\ne S_i\}\in\Lambda.
\]
These boxes form \(B_\Lambda\) and generate \(O_\Lambda\) on
\(\prod_i S_i\). Identify \(f,g\) when
\(\{i:f_i=g_i\}\in\Psi\), and write \(q(f)=[f]\).
The structure \(\prod^\Lambda_\Psi S_i\) has that quotient carrier,
coordinatewise functions, and basic relations true when their local truth
sets belong to \(\Psi\). Its topology \(O^\Lambda_\Psi\) is the
quotient topology from \(O_\Lambda\).
In its open sort the basic family is
\(B^\Lambda_\Psi=\{q(B):B\in B_\Lambda\}\).
This source's ordinary basis need not already contain the empty set.
Adjoin it when passing to the earlier basoid convention. A guarded
quantifier never has an empty-set witness, and its universal guard has
false antecedent there, so this addition changes no \(L_t\)-truth.

For a basic box \(U\), its canonical coordinate sets are
\(U_i=\pi_i[U]\), where \(\pi_i\) is the coordinate projection.
For a nonempty box these recover its ordinary factor opens. For an empty
box all canonical coordinate sets are empty. This convention is essential
for evaluating factor formulas with free open parameters: arbitrary written
factors of an empty Cartesian rectangle need not give its projections.
\end{hbdefinition}

\begin{hbtheorem}{thm:p03-topological-rip-structure}{proved=annex:p03-topological-rip-structure}{The open quotient map and its basic opens}
The point quotient is well defined, \(B_\Lambda\) is a basis, \(q\)
is an open quotient map, and \(B^\Lambda_\Psi\) is a basis for its
quotient topology. With canonical coordinates of a basic box,
\[
                 [f]\in q(U)
              \quad\Longleftrightarrow\quad
                 \{i:f_i\in U_i\}\in\Psi.
\]
For \(\Lambda=[I]^{<\omega}\), \(\Psi=\{I\}\), this is the
ordinary product topology; for \(\Lambda=\mathcal P(I)\),
\(\Psi=\{I\}\), it is the full box topology. The two support
conditions are distinct even though their underlying Cartesian structures
are the same.
\end{hbtheorem}

\begin{hbdefinition}{def:p03-topological-boolean-ideal}{The source's Boolean algebra with a distinguished ideal}
For index subsets put
\[
 A\equiv B\quad\Longleftrightarrow\quad
           \exists F\in\Psi\ (A\cap F=B\cap F).
\]
Here \(\equiv\) is the source's Boolean quotient congruence, not
elementary equivalence of structures. Write \([A]\) for its class and
use the source operations \([A]+[B]=[A\cup B]\),
\([A]\cdot[B]=[A\cap B]\), \([A]'=[I\setminus A]\),
\(0=[\varnothing]\), \(1=[I]\). The expanded structure is
\[
 \mathbf B=\langle\mathcal P(I)/{\equiv},+,\cdot,',0,1,
                                  \Lambda/{\equiv}\rangle,
 \qquad \Lambda/{\equiv}=\{[L]:L\in\Lambda\}.
\]
In \(L_B=\{+,\cdot,',0,1,\lambda\}\), \(\lambda\) names this
distinguished ideal. Its theory \(T_B\) consists of the Boolean-algebra
axioms together with
\(\lambda(0)\),
\(\lambda(u)\land\lambda(v)\to\lambda(u+v)\), and
\(\lambda(u)\land v\le u\to\lambda(v)\), where
\(v\le u\) abbreviates \(v\cdot u=v\).
Write \(u\setminus v=u\cdot v'\). These are Boolean terms and
first-order ideal axioms, not assertions of arbitrary join closure.
For a factor formula \(\psi\) and product/basic-box parameters set
\[
        I_\psi=\{i:(S_i,O_i)\models\psi[f_i,U_i]\}.
\]
Its Boolean class is distinct from the point class \([f]\), although
the source uses brackets for both; the domains declare their types.
\end{hbdefinition}

\begin{hblemma}{thm:p03-topological-carrier-substitution}{proved=annex:p03-topological-carrier-substitution}{Replacing an open parameter by the whole carrier}
For an \(L_t\)-formula in its membership presentation, let
\(\varphi^Y\) be obtained by replacing each free membership atom
\(t\varepsilon Y\) by \(t=t\), leaving bound occurrences unchanged
and using fresh bound variables. Then
\[
             (S,O)\models\varphi[v(Y/S)]
               \quad\Longleftrightarrow\quad
             (S,O)\models\varphi^Y[v].
\]
The substitution removes \(Y\) from the free variables and retains
the polarity of every other free open variable.
\end{hblemma}

\begin{hbdefinition}{def:p03-topological-fv-recursion}{Monotone Boolean reduction with support information}
A reduction of \(\varphi\) is a finite nonempty list of factor
\(L_t\)-formulas \(\psi_1,\ldots,\psi_m\) and a Boolean formula
\(\sigma\) increasing in all its displayed arguments, provably in
\(T_B\). Each factor formula has free variables among those of the
input and retains every consistently positive or negative free-open role
of the input. Use a truth component if a list would otherwise be empty.
An atom \(\alpha\) has reduction \((y=1;\alpha)\);
truth has \((y=1;\top)\) and falsity \((0\ne0;\top)\).
Conjunction concatenates the lists and conjoins the disjointly renamed
Boolean formulas. If \((\tau;\theta_1,\ldots,\theta_m)\) reduces
\(\chi\), negation uses
\[
                 (\neg\tau(y_1',\ldots,y_m');
                                   \neg\theta_1,\ldots,\neg\theta_m).
\]
This logical negation of the complemented-input condition is increasing;
it is not the decreasing formula \(\neg\tau(y_1,\ldots,y_m)\).

For either existential clause enumerate all subsets
\(s_1,\ldots,s_l\) of \(\{1,\ldots,m\}\), with the singletons
first and \(l=2^m\). Empty conjunction is truth. For an individual
quantifier \(\exists x\chi\), put
\[
 \psi_k=\exists x\bigwedge_{j\in s_k}\theta_j,\qquad
 \sigma(\bar y)=\exists z_1\cdots z_l
   \left[\bigwedge_k z_k\le y_k\land
    \bigwedge_{s_i\cup s_j=s_k}z_i\cdot z_j=z_k
                       \land\tau(z_1,\ldots,z_m)\right].
\]
For an open guard \(\exists Y(t\varepsilon Y\land\chi)\), with
\(Y\) negative in \(\chi\), use instead
\[
 \psi_k=\exists Y\left(t\varepsilon Y\land
                               \bigwedge_{j\in s_k}\theta_j\right),
 \qquad \eta_j=\theta_j^Y,
\]
and the source's condition
\[
 \sigma(\bar y,\bar v)=\exists z_1\cdots z_l
 \left[\bigwedge_k z_k\le y_k\land
   \bigwedge_{s_i\cup s_j=s_k}z_i\cdot z_j=z_k
   \land\tau(z_1,\ldots,z_m)
   \land\bigwedge_{j=1}^m\lambda(z_j\setminus v_j)\right].
\]
The local list is \((\psi_1,\ldots,\psi_l,\eta_1,\ldots,\eta_m)\).
Only the last ideal conditions constrain how many coordinates need a
proper neighborhood rather than the whole carrier. Boolean binders and
point/open binders are renamed capture-avoidantly before each clause.
\end{hbdefinition}

\begin{hbtheorem}{thm:p03-topological-fv}{proved=annex:p03-topological-fv}{Topological reduced-ideal-product composition}
For each fixed formula the recursion gives the stated increasing Boolean
condition and polarity-preserving local list, with
\[
 \prod^\Lambda_\Psi S_i\models\varphi[[f],q(U)]
          \quad\Longleftrightarrow\quad
             \mathbf B\models\sigma([I_{\psi_1}],\ldots,[I_{\psi_m}]).
\]
Here each free open parameter \(U\) is an actual basic box and its factor
values are its canonical coordinate projections. Satisfaction may use
either the quotient basis or its generated topology. Both existential
steps supply actual witnesses; the open witness is a box supported on a
set in \(\Lambda\). This is not merely the ordinary reduced-product
theorem: the distinguished ideal \(\lambda\) is indispensable to this
topological construction.
\end{hbtheorem}

\begin{hbcorollary}{thm:p03-topological-transfer}{proved=annex:p03-topological-transfer}{Theory preservation and topological ultraproducts}
For fixed \(I,\Lambda,\Psi\), factorwise \(L_t\)-equivalent
topological structures have \(L_t\)-equivalent reduced ideal-products.
For an ultrafilter \(\mathcal U\), take
\(\Lambda=\mathcal P(I)\) and \(\Psi=\mathcal U\).
Then for every \(L_t\)-formula with the stated point/basic-open
parameters,
\[
       \prod^{\mathcal P(I)}_{\mathcal U}S_i\models\varphi[[f],q(U)]
          \quad\Longleftrightarrow\quad I_\varphi\in\mathcal U.
\]
Its topology is generated by the quotient-box basis; it need not consist
of that basis alone. In a topological ultrapower the constant diagonal
preserves and reflects every \(L_t\)-formula on transported parameters.
Predicate interpretations use local truth sets that are ultrafilter-large,
including when a predicate is empty in some factor. They are not defined
as the image of an empty Cartesian product of predicate extensions.
\end{hbcorollary}
